% -*- coding: utf-8 -*-
\documentclass[a4paper,12pt]{article}
\usepackage{amsmath,amssymb}
\usepackage[russian]{babel}
\usepackage[utf8]{inputenc}
\usepackage[T2A]{fontenc}
\usepackage{geometry}
\geometry{margin=2cm}

\title{\Large \textbf{КОНТРОЛЬНЫЕ ВОПРОСЫ К ЗАЩИТЕ ЛАБОРАТОРНОЙ РАБОТЫ №2} \\
    \normalsize Численное решение нелинейных уравнений и систем}
\author{Вариант 2}
\date{2026}

\begin{document}

\maketitle

\section*{Ответы на контрольные вопросы}

\subsection*{1. Понятие точного и приближенного решений нелинейного уравнения}

\textbf{Ответ:} 

\textbf{Точное решение} (корень уравнения $f(x) = 0$) --- это значение $x^* \in \mathbb{R}$, такое что $f(x^*) = 0$ точно.

\textbf{Приближенное решение} (приближённый корень) --- это значение $x_n$, полученное численным методом, такое что:
\begin{itemize}
    \item $|f(x_n)| \leq \varepsilon$ (критерий по функции) или
    \item $|x_n - x^*| \leq \varepsilon$ (критерий по аргументу)
\end{itemize}

где $\varepsilon$ --- требуемая точность. Приближённое решение сходится к точному при $\varepsilon \to 0$.

\subsection*{2. Основная идея метода половинного деления?}

\textbf{Ответ:} 

На каждой итерации интервал $[a_n, b_n]$, содержащий корень (где $f(a_n) \cdot f(b_n) < 0$), делится пополам точкой $x_n = \frac{a_n + b_n}{2}$. Затем выбирается та половина отрезка, на концах которой функция имеет разные знаки, и процесс повторяется. Таким образом, интервал сужается на каждой итерации в два раза.

\textbf{Преимущество:} Гарантированная сходимость при единственном корне в интервале.

\subsection*{3. Может ли метод половинного деления найти точное значение корня уравнения?}

\textbf{Ответ:} 

Теоретически НЕТ, метод находит только приближённое решение. Практически, если корень лежит точно в одной из граничных точек интервала (например, $f(a_n) = 0$), то метод найдёт его на первой же итерации. Однако в общем случае для полиномов степени $n \geq 3$ и трансцендентных функций метод сходится к корню асимптотически, требуя бесконечного числа итераций для достижения точного значения.

\subsection*{4. В чем суть метода Ньютона?}

\textbf{Ответ:} 

Суть метода заключается в замене исходной функции $f(x)$ её касательной в точке текущего приближения $x_n$. Новое приближение находится как точка пересечения этой касательной с осью абсцисс.

\textbf{Рабочая формула:}
\[
    x_{n+1} = x_n - \frac{f(x_n)}{f'(x_n)}
\]

\textbf{Геометрический смысл:} На графике функции проводится касательная в точке $(x_n, f(x_n))$. Касательная пересекает ось $x$ в точке $x_{n+1}$.

\textbf{Скорость сходимости:} Квадратичная (очень быстрая).

\subsection*{5. Как выбирается начальное приближение для метода Ньютона?}

\textbf{Ответ:} 

Начальное приближение $x_0$ выбирается из условия \textbf{скорой сходимости}:
\[
    f(x_0) \cdot f''(x_0) > 0
\]

То есть, выбирается \textbf{тот конец интервала $[a, b]$}, где знак функции совпадает со знаком второй производной.

\textbf{Правило:} Если функция возрастает и выпукла вниз, начинаем с левого конца. Если убывает и выпукла вверх --- также с левого конца.

\textbf{Важно:} Неудачный выбор $x_0$ может привести к расходимости метода или сходимости к другому корню.

\subsection*{6. Идея метода хорд?}

\textbf{Ответ:} 

Метод хорд заменяет исходную функцию $f(x)$ на отрезке $[a, b]$ прямой линией (хордой), проходящей через точки $(a, f(a))$ и $(b, f(b))$. Новое приближение $x_n$ --- это точка пересечения этой хорды с осью абсцисс.

\textbf{Геометрический смысл:} Строится секущая через две точки на графике функции.

\textbf{Рабочая формула:}
\[
    x_n = a - \frac{b - a}{f(b) - f(a)} \cdot f(a) = \frac{a \cdot f(b) - b \cdot f(a)}{f(b) - f(a)}
\]

\textbf{Скорость сходимости:} Линейная (медленнее, чем Ньютон, но быстрее дихотомии).

\subsection*{7. Как выбирается начальное приближение для метода хорд с фиксированным концом интервала изоляции корня?}

\textbf{Ответ:} 

При использовании метода хорд с \textbf{фиксированным концом}, один конец интервала остаётся неподвижным на всех итерациях. Выбор фиксированного конца определяется правилом:

\textbf{Фиксируется тот конец, где $f(x) \cdot f''(x) > 0$,}

то есть где знак функции совпадает со знаком второй производной.

\textbf{Другой конец} служит подвижной точкой и инициализируется вторым концом интервала.

\textbf{Пример:} Если $f(a) \cdot f''(a) > 0$, то $a$ --- фиксирован, начальное приближение $x_0 = b$.

\subsection*{8. По каким причинам методы хорд и касательных предпочтительнее метода простой итерации?}

\textbf{Ответ:} 

\begin{enumerate}
    \item \textbf{Скорость сходимости}: 
    \begin{itemize}
        \item Метод Ньютона (касательных) --- \textbf{квадратичная} сходимость
        \item Метод хорд --- \textbf{линейная} сходимость (порядок $> 1$)
        \item Метод простой итерации --- \textbf{линейная} сходимость (медленнее)
    \end{itemize}
    
    \item \textbf{Требования к исходной информации}:
    \begin{itemize}
        \item Методы хорд и касательных требуют только интервал изоляции или одного начального приближения
        \item Метод простой итерации требует правильного преобразования уравнения $x = \varphi(x)$, что не всегда просто
    \end{itemize}
    
    \item \textbf{Универсальность}: Методы хорд и касательных применимы к широкому классу функций, методу простой итерации требуется специальное преобразование
    
    \item \textbf{Сходимость}: Методы хорд и касательных гарантируют сходимость при выполнении условий, метод простой итерации может не сходиться даже при выполнении условия $|\varphi'(x)| < 1$
\end{enumerate}

\subsection*{9. Какой из методов является двухшаговым методом? Как запустить этот метод?}

\textbf{Ответ:} 

\textbf{Метод секущих} является двухшаговым методом. Новое приближение определяется двумя предыдущими итерациями:
\[
    x_{n+1} = x_n - \frac{x_n - x_{n-1}}{f(x_n) - f(x_{n-1})} \cdot f(x_n)
\]

\textbf{Как запустить метод:}

Требуются два начальных приближения $x_0$ и $x_1$. Обычно выбирают:
\begin{itemize}
    \item $x_0$ и $x_1$ --- концы интервала изоляции $[a, b]$ или
    \item $x_0$ --- одно начальное приближение, $x_1 = x_0 + \delta$ (где $\delta$ --- малый шаг, например $\delta = 0.01$)
\end{itemize}

\subsection*{10. В чем суть метода простой итерации?}

\textbf{Ответ:} 

Метод простой итерации основан на \textbf{преобразовании} исходного уравнения $f(x) = 0$ к эквивалентному виду:
\[
    x = \varphi(x)
\]

Затем строится итерационная последовательность:
\[
    x_{n+1} = \varphi(x_n), \quad n = 0, 1, 2, \ldots
\]

которая при выполнении условий сходимости приближается к корню.

\textbf{Геометрический смысл:} На графике функции $y = \varphi(x)$ и прямой $y = x$ находятся точки пересечения, которые и являются корнями исходного уравнения. Итерационная последовательность представляет собой "ступеньки" или "спираль", сходящиеся к точке пересечения.

\subsection*{11. Каковы условия применяемости метода простой итерации?}

\textbf{Ответ:} 

Метод простой итерации применим, если выполняется:

\textbf{Достаточное условие сходимости:}
\[
    |\varphi'(x)| \leq q < 1 \quad \forall x \in [a, b]
\]

где $q$ --- коэффициент сжатия (константа Липшица).

\textbf{Скорость сходимости зависит от величины $q$:}
\begin{itemize}
    \item $q \approx 0$ --- быстрая сходимость
    \item $q \approx 1$ --- медленная сходимость
    \item $q \geq 1$ --- нет сходимости (расходимость)
\end{itemize}

\subsection*{12. Как правильно преобразовать исходное нелинейное уравнение $f(x) = 0$ к виду $x = \varphi(x)$?}

\textbf{Ответ:} 

Существует несколько способов:

\textbf{Способ 1: Прямое преобразование}
Выразить $x$ из одного из членов уравнения. Например:
\[
    f(x) = x^3 - x + 4 = 0 \Rightarrow x = x^3 + 4 \quad \text{(обычно не сходится)}
\]

\textbf{Способ 2: Методом с параметром (рекомендуемый)}
Вводится параметр $\lambda$ и преобразование:
\[
    x = x + \lambda f(x) \Rightarrow \varphi(x) = x + \lambda f(x)
\]

где $\lambda$ выбирается так, чтобы $|\varphi'(x)| < 1$.

Из условия оптимальности: $\varphi'(x) = 1 + \lambda f'(x) = 0 \Rightarrow \lambda = -\frac{1}{f'(x)}$

На практике берут: $\lambda = -\frac{1}{\max_{x \in [a,b]} |f'(x)|}$

\textbf{Способ 3: Переформулировка}
Иногда уравнение можно переписать, выделив отдельно, например:
\[
    \sin(x) + x = 1 \Rightarrow x = 1 - \sin(x)
\]

\textbf{Критерий правильности:} Функция $\varphi(x)$ должна удовлетворять условию $|\varphi'(x)| < 1$ на интервале $[a, b]$.

\subsection*{13. Каковы основные критерии окончания итерационного процесса?}

\textbf{Ответ:} 

Итерационный процесс останавливается при выполнении одного из следующих критериев:

\textbf{1. Критерий по аргументу:}
\[
    |x_n - x_{n-1}| \leq \varepsilon
\]
Применяется во всех методах.

\textbf{2. Критерий по функции:}
\[
    |f(x_n)| \leq \varepsilon
\]
Проверяет, насколько близко значение функции в найденной точке к нулю.

\textbf{3. Для метода половинного деления:}
\[
    |a_n - b_n| \leq \varepsilon
\]
Проверяет ширину оставшегося интервала.

\textbf{4. Для метода простой итерации (при $q > 0.5$):}
\[
    |x_n - x_{n-1}| < \frac{q}{1-q} \varepsilon
\]
Более строгий критерий для медленно сходящихся методов.

\textbf{5. Критерий по числу итераций:}
\[
    n > n_{\max}
\]
Предотвращает бесконечные циклы при расходимости.

\subsection*{14. Как оценить необходимое количество итераций в методе биссекции при заданной точности?}

\textbf{Ответ:} 

Для метода половинного деления интервал уменьшается вдвое на каждой итерации:
\[
    |a_n - b_n| = |a_0 - b_0| \cdot 2^{-n}
\]

Критерий остановки: $|a_n - b_n| \leq \varepsilon$

Отсюда:
\[
    |a_0 - b_0| \cdot 2^{-n} \leq \varepsilon
\]
\[
    2^{-n} \leq \frac{\varepsilon}{|a_0 - b_0|}
\]
\[
    2^{n} \geq \frac{|a_0 - b_0|}{\varepsilon}
\]
\[
    n \geq \log_2 \left( \frac{|a_0 - b_0|}{\varepsilon} \right)
\]

\textbf{Формула:}
\[
    n = \left\lceil \log_2 \left( \frac{b - a}{\varepsilon} \right) \right\rceil
\]

\textbf{Пример:} Для интервала $[0, 2]$ и точности $\varepsilon = 10^{-6}$:
\[
    n = \log_2(2 \cdot 10^6) \approx 20.93 \Rightarrow n = 21 \text{ итерация}
\]

\subsection*{15. Алгоритм решения системы нелинейных уравнений методом Ньютона?}

\textbf{Ответ:} 

\textbf{Система:}
\[
\begin{cases}
f(x, y) = 0 \\
g(x, y) = 0
\end{cases}
\]

\textbf{Алгоритм:}

\begin{enumerate}
    \item Выбрать начальные приближения $x_0, y_0$
    
    \item Вычислить матрицу Якоби (матрицу частных производных):
    \[
        J(x, y) = \begin{pmatrix} \frac{\partial f}{\partial x} & \frac{\partial f}{\partial y} \\ \frac{\partial g}{\partial x} & \frac{\partial g}{\partial y} \end{pmatrix}_{x=x_n, y=y_n}
    \]
    
    \item Вычислить определитель якобиана:
    \[
        \det J = \frac{\partial f}{\partial x} \cdot \frac{\partial g}{\partial y} - \frac{\partial f}{\partial y} \cdot \frac{\partial g}{\partial x}
    \]
    
    \item Решить систему линейных уравнений для приращений $\Delta x, \Delta y$:
    \[
        J \begin{pmatrix} \Delta x \\ \Delta y \end{pmatrix} = -\begin{pmatrix} f(x_n, y_n) \\ g(x_n, y_n) \end{pmatrix}
    \]
    
    По формулам Крамера:
    \[
        \Delta x = \frac{-f \cdot \frac{\partial g}{\partial y} + g \cdot \frac{\partial f}{\partial y}}{\det J}, \quad
        \Delta y = \frac{-g \cdot \frac{\partial f}{\partial x} + f \cdot \frac{\partial g}{\partial x}}{\det J}
    \]
    
    \item Обновить приближение:
    \[
        x_{n+1} = x_n + \Delta x, \quad y_{n+1} = y_n + \Delta y
    \]
    
    \item Проверить условие окончания: $\max(|\Delta x|, |\Delta y|) \leq \varepsilon$
    
    \item Если не выполнено, повторить с шага 2 для $n := n+1$
\end{enumerate}

\subsection*{16. Каковы преимущества и недостатки графического метода отделения решения для системы двух нелинейных уравнений?}

\textbf{Ответ:} 

\textbf{Преимущества:}
\begin{itemize}
    \item Наглядность и простота визуального поиска решений
    \item Позволяет определить количество решений системы
    \item Помогает выбрать хорошее начальное приближение для численных методов
    \item Не требует вычисления производных
\end{itemize}

\textbf{Недостатки:}
\begin{itemize}
    \item Неточность (зависит от качества графика)
    \item Затруднён в 3D и выше (система 3+ уравнений)
    \item Требует построения графиков обеих кривых (трудоёмко вручную)
    \item Может пропустить близко расположенные решения
    \item Не даёт точные координаты решений
\end{itemize}

\subsection*{17. В каких случаях можно применить метод простой итерации для решения системы нелинейных уравнений?}

\textbf{Ответ:} 

Метод простой итерации для систем применим при условии:

\[
\begin{cases}
x = \varphi_1(x, y) \\
y = \varphi_2(x, y)
\end{cases}
\]

\textbf{Достаточное условие сходимости (условие Липшица):}
\[
    \left| \frac{\partial \varphi_1}{\partial x} \right| + \left| \frac{\partial \varphi_1}{\partial y} \right| < 1
\]
\[
    \left| \frac{\partial \varphi_2}{\partial x} \right| + \left| \frac{\partial \varphi_2}{\partial y} \right| < 1
\]

\textbf{Практика:} Метод применяется, когда:
\begin{itemize}
    \item Систему удаётся преобразовать к нужному виду
    \item Выполняются условия сходимости
    \item Требуется простая реализация (не нужны производные)
    \item Метод Ньютона неприменим или сложен для реализации
\end{itemize}

\subsection*{18. Когда можно считать итерационный процесс законченным при использовании метода простой итерации для решения системы нелинейных уравнений?}

\textbf{Ответ:} 

Итерационный процесс считается законченным при выполнении одного из следующих условий:

\textbf{1. Критерий по аргументу:}
\[
    \max(|x_n - x_{n-1}|, |y_n - y_{n-1}|) \leq \varepsilon
\]

\textbf{2. Критерий по функции:}
\[
    \max(|f(x_n, y_n)|, |g(x_n, y_n)|) \leq \varepsilon
\]

\textbf{3. Взвешенный критерий (для медленно сходящихся методов при $q \approx 1$):}
\[
    \max(|x_n - x_{n-1}|, |y_n - y_{n-1}|) < \frac{q}{1-q} \varepsilon
\]

где $q = \max\left( \left| \frac{\partial \varphi_1}{\partial x} \right| + \left| \frac{\partial \varphi_1}{\partial y} \right|,  \left| \frac{\partial \varphi_2}{\partial x} \right| + \left| \frac{\partial \varphi_2}{\partial y} \right| \right)$

\textbf{4. Критерий по числу итераций:}
\[
    n \geq n_{\max}
\]

\subsection*{19. Что такое сходимость и скорость сходимости численных методов?}

\textbf{Ответ:} 

\textbf{Сходимость:} Числовой метод сходится, если итерационная последовательность приближений $\{x_n\}$ имеет предел, совпадающий с точным решением:
\[
    \lim_{n \to \infty} x_n = x^*
\]

\textbf{Скорость сходимости:} Характеризует, как быстро последовательность приближается к точному решению. Определяется порядком сходимости $p$:
\[
    |x_{n+1} - x^*| \leq C |x_n - x^*|^p
\]

\textbf{Классификация по порядку:}
\begin{itemize}
    \item $p = 1$ --- \textbf{линейная} сходимость (медленная). Примеры: половинное деление, простая итерация
    \item $1 < p < 2$ --- \textbf{сверхлинейная} сходимость (среднее). Пример: метод секущих ($p \approx 1.618$)
    \item $p = 2$ --- \textbf{квадратичная} сходимость (быстрая). Пример: метод Ньютона
    \item $p > 2$ --- \textbf{кубическая и выше} (очень быстрая)
\end{itemize}

\textbf{Практический смысл:} При квадратичной сходимости число верных цифр удваивается на каждой итерации, при линейной растёт медленнее.

\subsection*{20. Дайте определение устойчивости итерационного метода}

\textbf{Ответ:} 

\textbf{Устойчивость итерационного метода:} Метод называется устойчивым, если малые возмущения (ошибки) в исходных данных и ошибки округления приводят к малым ошибкам в конечном результате. Формально:

Метод устойчив, если при малых возмущениях $\delta x_0$ в начальном приближении и $\delta f$ в функции, погрешность в решении $\delta x^*$ остаётся малой и не накапливается по мере итераций.

\textbf{Математически:} Для метода $x_{n+1} = x_n + \lambda f(x_n)$ метод устойчив, если:
\[
    |\lambda f'(x)| \leq q < 1 \quad \Rightarrow \quad \text{ошибки затухают}
\]

\textbf{Примеры устойчивых методов:}
\begin{itemize}
    \item Половинное деление --- абсолютно устойчив
    \item Ньютон (при хорошем начальном приближении) --- устойчив локально
    \item Простая итерация (при $|\varphi'(x)| < 1$) --- устойчива
\end{itemize}

\textbf{Примеры неустойчивых методов:}
\begin{itemize}
    \item Ньютон (при плохом начальном приближении) --- может быть неустойчив
    \item Простая итерация (при $|\varphi'(x)| > 1$) --- неустойчива
\end{itemize}

\subsection*{21. Какой метод решения нелинейных уравнений наиболее чувствителен к выбору начального приближения (с точки зрения скорости сходимости)?}

\textbf{Ответ:} 

\textbf{Метод Ньютона (касательных)} наиболее чувствителен к выбору начального приближения.

\textbf{Обоснование:}
\begin{itemize}
    \item При $x_0$ близко к корню --- квадратичная сходимость (очень быстро)
    \item При $x_0$ далеко от корня --- может расходиться или сходиться к другому корню
    \item Если $f'(x_0) \approx 0$ или $f''(x_0) = 0$ --- метод может дать плохие результаты
\end{itemize}

\textbf{Рекомендация:} При применении метода Ньютона следует:
\begin{itemize}
    \item Предварительно отделить корни графически или табулированием
    \item Выбрать $x_0$ из условия $f(x_0) \cdot f''(x_0) > 0$
    \item Иногда применить сначала несколько шагов половинного деления для грубого приближения
\end{itemize}

\subsection*{22. Возможно ли применение метода Ньютона для решения уравнения $x^3 - x^2 - 2.5x + 2 = 0$ на интервале [5, 6]?}

\textbf{Ответ:} 

Сначала проверим, содержит ли интервал [5, 6] корень:

$f(5) = 125 - 25 - 12.5 + 2 = 89.5 > 0$

$f(6) = 216 - 36 - 15 + 2 = 167 > 0$

Поскольку $f(5) \cdot f(6) > 0$ (одинаковые знаки), на интервале [5, 6] \textbf{нет корня}.

\textbf{Вывод:} Применить метод Ньютона на интервале [5, 6] \textbf{НЕЛЬЗЯ}, так как там нет корня.

Рекомендуется:
\begin{itemize}
    \item Расширить поиск на другие интервалы
    \item Построить график функции для поиска корней
    \item Использовать табулирование функции
\end{itemize}

\subsection*{23. Пусть применен один шаг метода хорд для решения нелинейного уравнения $x^3 - x^2 - 2.5x + 2 = 0$ на интервале [5, 6]. Какой интервал будет получен для дальнейшего вычисления корня?}

\textbf{Ответ:} 

Из ответа на вопрос 22: функция имеет одинаковые знаки на концах интервала [5, 6], поэтому на этом интервале нет корня, и \textbf{метод хорд не применим}.

Однако, если предположить гипотетический случай, где интервал содержит корень, то новый интервал выбирается следующим образом:

После первого шага метода хорд получаем точку $x_1$. Вычисляем $f(x_1)$ и выбираем новый интервал:
\[
[a_1, b_1] = \begin{cases}
[a_0, x_1], & \text{если } f(a_0) \cdot f(x_1) < 0 \\
[x_1, b_0], & \text{если } f(x_1) \cdot f(b_0) < 0
\end{cases}
\]

То есть, выбирается та часть интервала, где функция меняет знак.

\subsection*{24. За какое количество итераций возможно решение нелинейных уравнений методом бисекций на отрезке [0; 2] с точностью $10^{-3}$?}

\textbf{Ответ:} 

Используем формулу из ответа на вопрос 14:
\[
    n = \left\lceil \log_2 \left( \frac{b - a}{\varepsilon} \right) \right\rceil
\]

Где:
\begin{itemize}
    \item $a = 0, b = 2$, значит $b - a = 2$
    \item $\varepsilon = 10^{-3} = 0.001$
\end{itemize}

\[
    n = \log_2 \left( \frac{2}{0.001} \right) = \log_2(2000) \approx 10.97
\]

\textbf{Ответ:} $n = 11$ итераций

\subsection*{25. Сколько итераций необходимо для решения уравнения $x^2 + 0 = 0$ методом бисекций с точностью $\varepsilon = 10^{-6}$ на отрезке [-3; -1]?}

\textbf{Ответ:} 

\textbf{Замечание:} Уравнение $x^2 + 0 = 0$ эквивалентно $x^2 = 0$, которое имеет единственный корень $x^* = 0$.

Однако этот корень \textbf{НЕ находится} на интервале [-3, -1], так как:
\begin{itemize}
    \item $f(-3) = 9 > 0$
    \item $f(-1) = 1 > 0$
\end{itemize}

Поскольку $f(-3) \cdot f(-1) > 0$ (одинаковые знаки), на интервале [-3, -1] \textbf{нет корня}.

\textbf{Вывод:} Метод половинного деления \textbf{НЕПРИМЕНИМ} на интервале [-3, -1], так как там нет корня.

Корректный интервал для этого уравнения должен содержать 0, например [-1, 1].

На интервале [-1, 1]:
\[
    n = \log_2 \left( \frac{1 - (-1)}{10^{-6}} \right) = \log_2(2 \cdot 10^6) \approx 20.93
\]

\textbf{Ответ:} На интервале [-1, 1] требуется $n = 21$ итерация.

\end{document}
